\section{Related Work}

\subsection{Image-Based Plant Disease Classification}

Image-based plant disease classification has a substantial literature,
particularly since the public release of large curated crop image
collections such as PlantVillage~\citep{hughes2015plantvillage}.
\citet{mohanty2016plantdisease} demonstrated that deep convolutional
networks pretrained on ImageNet~\citep{russakovsky2015imagenet} can be
transferred to plant disease classification with strong in-distribution
performance, establishing the basic feasibility of smartphone-style
diagnosis. Most subsequent work in this tradition terminates at the class
label and does not address downstream questions of explanation,
uncertainty quantification, or evidence selection. AGAI inherits the
classifier-as-anchor design from this body of work but treats explanation
as a separate stage with its own constraints.

ResNet-50~\citep{he2016resnet} remains a strong choice for this setting
because residual connections support deep feature extraction without the
optimization instability that historically affected very deep CNNs. We do
not claim novelty for the classifier architecture; rather, the contribution
of the classifier in our pipeline is that it is calibrated, gated, and
isolated from the language model.

\subsection{Probability Calibration}

Modern neural classifiers are commonly miscalibrated, producing
overconfident maximum-softmax probabilities. \citet{guo2017calibration}
showed that a single-parameter post-hoc temperature scaling correction
substantially reduces calibration error without affecting accuracy. We
adopt this recipe. Expected calibration error
(ECE)~\citep{naeini2015ece} is reported for the deployed classifier, and
Wilson score intervals~\citep{wilson1927interval} are used for per-class
accuracy uncertainty due to the small per-class support of the balanced
holdout.

\subsection{Detection Transformers and Visible-Part Grounding}

The DETR family~\citep{carion2020detr} formulates object detection as a
direct set prediction problem, eliminating hand-tuned anchors and
non-maximum suppression. RF-DETR~\citep{robinson2025rfdetr} is a
real-time variant in this family and forms the basis of our plant-part
detector. In AGAI, the detector serves a non-standard purpose: rather
than localize disease symptoms, it produces a binary mask over a fixed
six-part vocabulary that is injected into the language model prompt as a
visibility constraint. This converts a downstream-only inference task
(``what plant parts may be discussed?'') into an upstream perception
task with an explicit ground truth.

\subsection{Vision-Language Models and Parameter-Efficient Adaptation}

The MiniGPT line~\citep{zhu2023minigpt4,chen2023minigptv2} aligns a
vision encoder with a large language backbone via a learned linear
projection and lightweight instruction tuning. The closely related LLaVA
family~\citep{liu2023llava} demonstrates the broader utility of this
class of architectures. AGAI uses MiniGPT-v2 with a Llama 2 chat
backbone~\citep{touvron2023llama2} adapted via low-rank
adaptation~\citep{hu2021lora}. The architectural choices are pragmatic;
the contribution lies in how the language model is wrapped: as a
constrained report writer downstream of an authoritative classifier and
detector, rather than as a free diagnostician.

\subsection{Retrieval-Augmented Generation and Grounded Output}

Retrieval-augmented generation~\citep{lewis2020rag} attaches an external
non-parametric memory to a generative model so that factual content can
be sourced from a curated corpus rather than parameterized weights. We
follow this pattern at runtime: a curated disease knowledge graph,
populated from cooperative-extension publications, supplies optional
treatment-context snippets to the report-writer prompt. The retrieval
stage is not the diagnostic authority; it provides citation-backed
context for the explanation.

\subsection{Hallucination in Multimodal Models}

The recent survey of \citet{huang2024hallucinationsurvey} catalogues
the landscape of hallucination behaviors in large language models and
multimodal systems, including object hallucination, attribute
hallucination, and unsupported causal claims. AGAI's design responds to
this literature by reducing the surface area on which the language model
can hallucinate: the classifier label is fixed, the visible-part list is
fixed, and the generation is restricted to explanation conditioned on
both. The constrained dual-report supervision policy described in
\S\ref{sec:system} extends this idea to training-data generation.

\subsection{Locality and Inference-Time Cost}

Several recent agricultural-AI deployments rely on hosted multimodal APIs
at inference time. AGAI deviates from this pattern: the deployed runtime
is fully local, with no external model calls during serving. GPT-class
models are used only offline during stage-2 supervision generation. Any
quality improvement therefore must survive deployment without privileged
runtime access to a frontier model.
